\section{Extensions of our result}

In this section, we illustrate the modularity of our proof method by showing that it extends to more complicated systems.
In fact, we believe that the proof should extend to various flavours of type theories, as long as they satisfy definitional inversion.

For instance, we describe how to extend our equivalence result to the theory of \citet{polymorphism, bcde2026}, with explicit universe polymorphism and cumulative universes indexed by internal levels.
For the theory without constraints, the proof of the previous sections can be adapted directly by simply making the reasoning on levels internal to the theory.
However, for a theory with constraints, things get more complicated. 

In this system, contexts can be extended with level variables $\Gamma, \alpha$ and level constraints $\Gamma, \psi$, we also have level and constraint products $[\alpha]A$, $[\psi]A$, along with a collapse element $\emptyset$ in inconsistent level contexts.

First, to relate the Tarski and Russell presentations ($\mathcal{T}_T$ and $\mathcal{T}_R$). When we look closely at what was needed in Section \ref{sectionT1T2}, we observe that for universe levels we only need to take suprema and infima. In fact, the lower universe $\U_k$ of Lemma \ref{mainuniverse} is given by $\min(n_0,n_1)$ in the base cases. To reconstruct this operation internally, we observe that in any consistent context, if a constraint $l \leq n_0$ and $l \leq n_1$ holds, we can use the decidability of the arithmetic on levels.

Second, to relate the annotated and unannotated Russell presentations ($\mathcal{T}_R$ and $\mathcal{T}_P)$, where $\strip$ also erases the type annotations on $\langle \alpha \rangle_A u$, $\langle \psi \rangle_A t$, and $t_A l$), we extend the join and coercion method of Section \ref{sectionT2T2'}. Because cumulativity does not act structurally on $[\alpha]A$ or $[\psi]A$, we extend Definition \ref{def:join} with join rules for $[\alpha]B_i$ and $[\psi]B_i$, where universe joins use the internal supremum $\U_{k \vee m}$ and the corresponding coercions are defined by $\eta$-expansion.

\[
    \coe_{[\alpha]B \leq [\alpha]L'}(t) := \langle \alpha \rangle_{L'}\,\coe_{B \leq L'}((t\p)_{B\p^+}\,\alpha), \qquad
    \coe_{[\psi]B \leq [\psi]L'}(t) := \langle \psi \rangle_{L'}\,\coe_{B \leq L'}(t).
\]
For the diagonal lemma (Lemma \ref{lem:join-diag}) to go through for constraint products, we add the natural $\eta$-rule for constraint products :
\[\Gamma \vdash \langle \psi \rangle_A t  = t : [\psi]A \quad \text{when} \quad \Gamma \vdash t : [\psi]A\]
Moreover, since $[\psi]A = A$ holds whenever $\psi$ is valid in $\Gamma$, a constraint product is not necessarily distinct from a universe, $\Pi$-type, or level product. Thus, the domain model should have a guard-aware no-confusion property stating that in a consistent context, if $[\psi]A$ is convertible to $\U_l$, $\Pi(C,B)$, or $[\alpha]B$, then $\psi$ is valid in $\Gamma$.
